\section{The second conductor-sixteen Tate residue on the original field}
\label{sec:conductor16-second-residue-strict}

\paragraph{Status and exact dependency.}
The only class-group input is AC-G144-class, specifically
$\operatorname{Cl}(L)[269]=0$ for the field $L$ below.
Its provisional class-number selector is not used.
The strict proof and exact local computation are independent of
the analytic class-certificate computation.  Review status and
the exact class dependency are recorded separately.
No automorphic seed, auxiliary host, model, or mixed realization
is assumed.  Every coefficient ring in this section is the original one.

Put $p=269$, $K=\mathbb Q(\omega)$, and $t=\sqrt[3]{p}$ with its
real embedding.  Fix $\chi:G_K\longrightarrow\mu_3$ for $K(t)/K$.
Let $\phi$ be the arithmetic quartic character modulo sixteen with
\[
 \ker\phi=\{1,7\},\qquad \phi(13)=i,\qquad
 \phi(5)=-i,\qquad i\equiv82\pmod p .
 \tag{C16T.1}
\]
Then $\phi(-1)=-1$ and $\phi^2=\chi_8$.
Write $c_p$ for the actual cyclotomic character, $\omega_p$ for
its Teichmuller part, and
\[
             \tau_p=\omega_p^{134}.
 \tag{C16T.2}
\]
The last character is the quadratic character of
$\mathbb Q(\sqrt p)$: the quadratic Gauss-sum identity gives
the field $\mathbb Q(\sqrt{(-1)^{(p-1)/2}p})$, and $p\equiv1\bmod4$.
Arithmetic reciprocity is used throughout.

For a discrete coefficient $D$ define the original strict group by
\[
 \begin{gathered}
 \operatorname{Sel}_{\rm str}(K,D)=
 \ker\bigl(H^1_{\rm cts}(G_K,D)\longrightarrow\mathcal L(D)\bigr),\\
 \mathcal L(D)=
 \prod_{v\mid p}H^1(G_{K_v},D)\ \times\
 \prod_{v\nmid p,\infty}H^1(I_v,D).
 \end{gathered}
 \tag{C16T.3}
\]
Thus the condition at $p$ is zero in full local cohomology.

\begin{theorem}[AC-G144: conductor-sixteen strict vanishing at the second residue]
\label{AC-G144}
Let $E/\mathbb Q_p$ be finite, let $O$ be its original ring of integers,
and suppose that $O$ contains the values of $\chi$.
For
\[
 \begin{gathered}
 j\in\{1,2\},\quad \sigma\in\{1,-1\},\quad m\in\mathbb Z,\quad e\geq1,\\
 \kappa:G_K\longrightarrow1+\mathfrak m_O
       \quad\hbox{continuous and unramified away from }p,
 \end{gathered}
 \tag{C16T.4}
\]
put
\[
       D=(E/O)(\chi^{-j}\phi^{-\sigma}
                        \omega_p^{-m}\kappa^{-1}c_p^{-e}).
 \tag{C16T.5}
\]
If $e+m\equiv133\bmod268$, then
\[
                 \operatorname{Sel}_{\rm str}(K,D)=0 .
 \tag{C16T.6}
\]
The corresponding residual strict group also vanishes over every
finite residue field containing its character values.
There is no absolute ramification restriction on $O$.
In particular, for $m=0$ and $\kappa=1$, this proves strict vanishing
for the same four finite characters $\chi^{-j}\phi^{-\sigma}$ that
occur in the actual AC-K25 hosts, at all $e\equiv133\bmod268$.
The auxiliary quadratic character $\tau_p$ belongs to the residual
Kummer splitting field; it is not added to those four finite labels.
\end{theorem}

\begin{proof}
\proofstep{1. The exact twisted quartic field.}
Let $z=\zeta_{16}$, put $w_0=z-z^{-1}$, and choose
\[
 \begin{gathered}
 w=\sqrt p\,w_0,\qquad F=\mathbb Q(w),\\
 L=F(t),\qquad M=KL,\qquad
 w^4+4p w^2+2p^2=0 .
 \end{gathered}
 \tag{C16T.7}
\]
The fields $\mathbb Q(w_0)$ and $\mathbb Q(\sqrt p)$ are disjoint,
since the first is unramified at $p$.
In their product, the simultaneous sign change of $w_0$ and
$\sqrt p$ fixes $w$.  The resulting quotient is cyclic of order four.
It is exactly the quotient cut out by
\[
                         \psi=\phi\tau_p.
 \tag{C16T.8}
\]
Indeed the generator with $\phi=i$ and $\tau_p=1$ sends
$w$ to the next quartic conjugate, while the nontrivial quadratic
element has the same action as its square.  Its action is
\[
 b(w)=-3w-\frac{w^3}{p},\qquad b^2(w)=-w,\qquad \psi(b)=i .
 \tag{C16T.9}
\]
The unique quadratic subfield is $\mathbb Q(\sqrt2)$.
The field $F$ is unramified at three.
The $S_3$ splitting field $K(t)$ has only one nontrivial abelian
quotient, namely $K$, which is ramified at three.
Consequently $F\cap K(t)=\mathbb Q$, and
\[
 \operatorname{Gal}(M/\mathbb Q)=S_3\times C_4,\qquad
 \Delta:=\operatorname{Gal}(M/K)=C_3\times C_4.
 \tag{C16T.10}
\]
In particular $[L:\mathbb Q]=12$ and $\mu_p\not\subset M$,
since $p-1=268$ does not divide $[M:\mathbb Q]=24$.

The reduced generator used for the independent class calculation
is $-u$, where
\[
 \begin{gathered}
 r=\frac{w^2}{p},\qquad u=\frac wt,\qquad r^2+4r+2=0,\\
 u^6=p r^3,\qquad t=\frac{u^2}{r},\qquad w=\frac{u^3}{r}.
 \end{gathered}
 \tag{C16T.11}
\]
Since $r^3=14r+8$, the element $u$ determines
$r=(u^6/p-8)/14$, and then $t,w$.
Thus $L=\mathbb Q(u)$, with minimal polynomial
\[
                    u^{12}+40p\,u^6+8p^2.
 \tag{C16T.12}
\]
This checks the precise field used in AC-G144-class.

\proofstep{2. The two ramified completions.}
The polynomial $r^2+4r+2$ has discriminant eight.
Since $p\equiv5\bmod8$, it is irreducible modulo $p$.
Therefore $A_0=\mathbb Q_p(r)$ is the unramified quadratic
extension, which is also isomorphic to $K_p$.
The polynomial $u^6-pr^3$ is Eisenstein over $A_0$.
Equations (C16T.11) show that
$A_0(u)$ contains the required images of $K,t,w$.
It has degree twelve over $\mathbb Q_p$, ramification degree six,
and residue degree two.
Since $M/\mathbb Q$ is Galois of degree twenty-four, there are
exactly two places of $M$ above $p$, each with
\[
                             e_p=6,\qquad f_p=2.
 \tag{C16T.13}
\]
At one chosen place, $u$ is a uniformizer, and the residue field is
\[
                         \mathbb F_p[r]/(r^2+4r+2).
 \tag{C16T.14}
\]
The global actions satisfy
\[
 b(r)=-4-r,\qquad b(u)=-(3+r)u,\qquad b^2(u)=-u.
 \tag{C16T.15}
\]
The polynomial identity
$(3+r)(3+(-4-r))=-1$ verifies the last equality.
The element $b$ exchanges the two places: it acts nontrivially
on the residue of $r$, whereas it fixes $K$, whose residue field
at the chosen place already equals the full quadratic residue field.
The decomposition group inside $\Delta$ is
$C_3\times\langle b^2\rangle$.
In all that follows the target is the direct sum over both places.
No $K_p$-linear $b$-action on a single completion is asserted.

\proofstep{3. Inverse-character Kummer descent and the class sector.}
Let $k=O/\mathfrak m_O$.  All fourth roots of unity already lie
in $\mathbb Q_p$, so no coefficient extension is needed for $\phi$.
Reduction of the character in (C16T.5) gives
\[
 \begin{aligned}
 a&=\bar\chi^{-j}\bar\phi^{-\sigma}\bar c_p^{-(e+m)},\\
 \eta:=a/\bar c_p
   &=\bar\chi^{-j}\bar\phi^{-\sigma}\bar\tau_p
     =\bar\chi^{-j}\bar\psi^{-\sigma}.
 \end{aligned}
 \tag{C16T.16}
\]
Here $-(e+m+1)\equiv134\bmod268$ and
$\tau_p^{-1}=\tau_p^{-\sigma}=\tau_p$.
The character $\eta$ is faithful on $\Delta$ and has splitting
field $M$.  Prime-to-$p$ inflation--restriction, globally and
at every local subgroup, identifies the residual strict kernel with
\[
 e_{\eta^{-1}}\ker\left(
       V_{M,S}\otimes_{\mathbb F_p}k
           \longrightarrow\mathcal P_M\otimes_{\mathbb F_p}k\right),
 \tag{C16T.17}
\]
where
\[
 \begin{aligned}
 V_{M,S}&=\{x\in M^\times:p\mid v(x)\text{ for }v\nmid p\}/M^{\times p},\\
 V_M&=\{x\in M^\times:p\mid v(x)\text{ for every finite }v\}/M^{\times p},\\
 \mathcal P_M&=\bigoplus_{v\mid p}M_v^\times/M_v^{\times p}.
 \end{aligned}
 \tag{C16T.18}
\]
This uses ordinary Kummer theory on $M$: the coefficient restricts
to $k(1)$, with the extra $\Delta$-twist $\eta$ on its cohomology.
Taking invariants therefore selects $\eta^{-1}$, not $\eta$.
Away from $p$, the inertia-zero condition is exactly divisibility
of the Kummer valuation by $p$.  A local $p$th power at $p$ also
has valuation divisible by $p$, so the kernel in (C16T.17)
belongs to $V_M$.

There is an exact ordinary-unit sequence
\[
 0\longrightarrow \mathcal O_M^\times/\mathcal O_M^{\times p}
 \longrightarrow V_M\longrightarrow\operatorname{Cl}(M)[p]
 \longrightarrow0.
 \tag{C16T.19}
\]
Use AC-G144-class only through $\operatorname{Cl}(L)[p]=0$.
Let $s$ be the reflection in the $S_3$ factor fixing $L$.
Extension and norm have composites $2$ and $1+s$.
Thus $\operatorname{Cl}(M)[p]^s=0$, and $s$ acts by $-1$.
Every conjugate reflection acts by $-1$ as well.
Their products generate $C_3$, so $C_3$ acts trivially on
$\operatorname{Cl}(M)[p]$.  Its nontrivial cubic part is zero.
In particular no vanishing of the entire class group of $M$ is assumed.

\proofstep{4. The exact unit sector has dimension two.}
The twelve complex places of $M$ above the complex place of $K$
form a regular $\Delta$-set.  The equivariant Dirichlet unit
character is consequently $\mathbb Q[\Delta]-\mathbb Q$.
Since $p\nmid12$ and $\mu_p(M)=1$, reduction of the integral unit
module preserves these character multiplicities.
For each $\lambda=82,187$, the nontrivial cubic part of the
$b$-eigenspace with eigenvalue $\lambda$ is a two-dimensional
$\mathbb F_p$-space $W_\lambda$.
It is simple as an $\mathbb F_p[C_3]$-module because
$p\equiv2\bmod3$ and $X^2+X+1$ is irreducible.
Therefore one nonzero value of the equivariant localization
map on $W_\lambda$ proves injectivity on that whole space.
After scalar extension to $k$, both cubic eigenlines inject.

\proofstep{5. An exact global unit and all its denominator terms.}
Use unit number four in the exact coefficient file
\path{evidence/G144_sources/strict/unit4_exact.json}.
Its coefficient basis is $t^a w^b$, $0\leq a\leq2$,
$0\leq b\leq3$, ordered by $4a+b$.
The file records the twelve integer numerators, common denominator
$3p^2=217083$, and a monic integral degree-twelve relation with
constant coefficient one.  Its SHA256 is, split for printing,
\begin{quote}\ttfamily\small
1957e60e3b6f6e42dede3af88e2b3ece\\
12e0d0b54a0933ba39cc22c3d208128b.
\end{quote}
The rational vector specifies an exact element $U\in L$.
The monic relation proves integrality; its reversed relation proves
that $U^{-1}$ is integral.  Thus $U$ is an actual global unit.
Neither a provisional fundamental-unit basis nor a numerical
regulator is used for this conclusion.

The independent verifier converts this entire rational vector
to the exact local basis $1,u,\ldots,u^5$ over $\mathbb Q(r)$.
For every summand $c_{ab}t^a w^b$, write
$2a+3b=6q+j$ with $0\leq j<6$.  Its contribution is exactly
\[
             c_{ab}p^q r^{3q-a-b}u^j.
 \tag{C16T.20}
\]
All fractions, including the full $p^2$ denominator, are retained
until coefficients with the same power of $u$ have been combined.
For each of the six resulting coefficients, both rational
coordinates in the basis $1,r$ are $p$-integral.
Horner evaluation in
$\mathbb Q(r)[u]/(u^6-pr^3)$ independently verifies the monic
relation, without reusing a multiplication-matrix characteristic
polynomial.  Only after these exact checks does reduction give
\[
            U=A+uB\pmod{u^2},\qquad
            A=95+232r,\qquad B=38+21r.
 \tag{C16T.21}
\]
Here $A\ne0$ in (C16T.14).

Let $c$ generate the cubic factor, so $c(t)=\omega t$,
$c(w)=w$ and $c(u)=\omega^{-1}u$.
For the global unit
$V=U^3/(U\,c(U)\,c^2(U))$, the sum of the three linear
cubic terms is zero.  Hence
\[
              V=1+uR\pmod{u^2},\qquad
              R=3B/A=94+19r.
 \tag{C16T.22}
\]
For each $\lambda=82,187$, choose integers representing
$\lambda^{-a}\bmod p$ and form the global product
$\prod_{a=0}^3 b^a(V)^{n_a}$.
The four linear coefficients before weighting are
\[
 \begin{array}{c|c}
 a&\text{coefficient of }u\text{ in }b^a(V)\\ \hline
 0&94+19r\\
 1&177+232r\\
 2&175+250r\\
 3&92+37r .
 \end{array}
 \tag{C16T.23}
\]
For the odd powers this calculation uses
$b(r)=-4-r$ and $b(u)=-(3+r)u$ from (C16T.15).
The weighted results are
\[
 \begin{array}{c|c|c}
 \lambda&(n_0,n_1,n_2,n_3)\bmod p&
                    \text{coefficient of }u\\ \hline
 82&(1,187,268,82)&212+188r\\
 187&(1,82,268,187)&164+157r .
 \end{array}
 \tag{C16T.24}
\]
Both are nonzero in the quadratic residue field.
The product was formed globally before evaluating at the chosen
completion.  Thus (C16T.23)--(C16T.24) do not identify the two
completions or assume that $b$ acts locally on just one of them.

Modulo $p$th powers, $3-(1+c+c^2)$ kills the trivial cubic
part and is multiplication by three on its complement.
The operator $\sum_a\lambda^{-a}b^a$ is four times the
quartic idempotent.  Both omitted scalars are invertible.
Thus the constructed products belong to the images of
the two required spaces $W_\lambda$.
Each has residue one.  If it were a local $p$th power, its
root would also have residue one because Frobenius is an
automorphism of the residue field.  Write that root as
$1+uy$.  In $(1+uy)^p$ the linear term has valuation at
least $v_u(pu)=7$, and every other nonconstant term has
valuation at least two.  Its $u$ coefficient modulo $u^2$
would be zero, contradicting (C16T.24).
Localization is therefore nonzero, and hence injective,
on both simple sectors.  Together with (C16T.17)--(C16T.19)
this proves the residual strict vanishing.

\proofstep{6. The full lift on the original coefficient ring.}
Let $\varpi$ be a uniformizer of the given $O$.
The residual inertia of $\chi$ at $p$ has order three.
The residual cyclotomic factor has order dividing $268$,
$\phi$ is unramified there, and $\bar\kappa=1$.
They cannot cancel the nontrivial cubic inertia.
Some inertia element therefore acts on $D$ by a scalar whose
difference from one is a unit.  Consequently
$D^{G_{K_p}}=0$ and $D^{G_K}=0$.
Away from $p$, both $c_p$ and $\kappa$ are unramified.
The remaining finite inertia character has order dividing twelve,
prime to $p$.  Its invariants are zero if the character is nontrivial,
and otherwise are all of $D$, which is $\varpi$-divisible.
For every subgroup $J$ used in (C16T.3) this proves
\[
                         D^J/\varpi D^J=0.
 \tag{C16T.25}
\]
Apply the global and local long exact sequences of
$0\to D[\varpi]\to D\xrightarrow{\varpi}D\to0$.
A strict class killed by $\varpi$ lifts globally to
$H^1(K,D[\varpi])$.  Equation (C16T.25) makes the map
$H^1(J,D[\varpi])\to H^1(J,D)$ injective at every local subgroup,
so that lift is strict.  Residual strict vanishing makes it zero.
Finally every continuous cocycle from the compact group $G_K$
to the discrete module $D$ has finite image, and is killed by
some power of $\varpi$.  A nonzero strict class would yield
a nonzero strict class killed by $\varpi$, a contradiction.
This proves (C16T.6), without enlarging or replacing the original $O$.
\end{proof}

\paragraph{Actual-host scope and the distinction between labels and splitting fields.}
The character $\phi$ factors through
$\mathbb Q(\zeta_{2^{h+2}})$ for $h\geq2$, while $\chi$ factors
through the actual cubic part of
\[
 H_u=\mathbb Q(\zeta_{2^{h+2}},\zeta_7,\sqrt{-q},\sqrt{-d},\omega,t).
\]
Thus the four finite labels at $m=0$, $\kappa=1$ are actual
AC-K25 labels.  The field $F$ of (C16T.7) is used only to
compute their residual Kummer classes at the second Tate residue.
It is not asserted to lie in $H_u$, and no new host character
$\tau_p$ is introduced.
For a positive integer $T$ and $1\leq e<T$ this theorem supplies exactly
$4\lfloor(T+134)/268\rfloor$ additional strict-zero labels.
Combining it with AC-G142 gives
$4\lfloor T/134\rfloor$ strict-zero labels at
$e\equiv133\bmod134$.
Transferring those zero strict groups to positive determinant indices
still retains the local comparison terms of AC-K26-local,
as in AC-G143.  No individual index-zero assertion, remaining
all-prime estimate, mixed realization, or analytic zero exclusion
is inferred here.

\paragraph{Geometric dictionary and source scope.}
The proof concerns rank-one Artin--Tate sheaves on the arithmetic
curve of $K$, the finite cover split by $M/K$, the Kummer
unit/class sequence, and restriction to every local stalk.
The finite-cover/representation dictionary is the retained
SGA~1, Expos\'e~V, \S\S4 and 7 input of AC-K25;
the Kummer sequence is Stacks tag~03PK, as used in AC-G103,
AC-G130, AC-G134 and AC-G142.
Their earlier applied-source audits are retained; no new
whole-primary-source reading or Lean execution is claimed.
The new finite input is the exact ramified expansion and
quartic projection above.  The distinct class-number input is
AC-G144-class, with its own analytic and finite certificates.
